% GENERATED FILE. Edit canonical lessons, then run npm run generate:books.
% canonical-source-hash: fnv1a64:a165e5eb749efbc1
% canonical-lessons: PT-C408-peao, PT-C408-passadeira, PT-C408-atropelar, PT-C408-radar, PT-C408-fiscalizar, PT-R408-revisao-401-408

\chapter{Peões e radares}
\label{ch:pt-peoes-e-radares}
\hlchaptermodality{408}

\begin{chapteropening}
\textbf{By the end of this chapter:} \emph{I can talk about road rules in town: pedestrians, the zebra crossing, running someone over, speed cameras and road checks.}

Complete the last lesson of chapter 408: I can talk about road rules in town: pedestrians, the zebra crossing, running someone over, speed cameras and road checks.
\end{chapteropening}

\section[o peão]{o peão — the pedestrian; the pawn (in chess)}

\label{lesson:PT-C408-peao}



\begin{quote}
Before the new one: say the Portuguese for the screw, then the Portuguese for to unscrew.
\end{quote}

\subsection*{You'll want to know: o peão}
\textbf{o peão} — \textquotedblleft{}the pedestrian; the pawn (in chess)\textquotedblright{}.

\begin{quote}
Esta rua é só para peões.
\end{quote}

\begin{quote}
\emph{This street is only for pedestrians.}
\end{quote}

In Brazil: \emph{o pedestre}.

\begin{cousinweb}
\emph{peão} is inherited from Late Latin \emph{pedō, pedōnis} 'foot soldier', built on \emph{pēs} 'foot' (EtymDB), which also lists Spanish \emph{peón} and French \emph{pion} as its cousins.
\end{cousinweb}

\begin{grammarlens}[title={What you've built}]
One more word for telling, explaining and describing at B1.
\end{grammarlens}

\subsection*{Your turn}
\begin{itemize}
\raggedright
  \item \emph{Say it:} \emph{o peão}
  \item \emph{Say it:} \emph{o peão}, once more
  \item \emph{Say it:} a sentence of your own with \emph{o peão}
  \item \emph{Recall:} say the Portuguese for the screw, then the Portuguese for to unscrew, then say \emph{o peão} again
\end{itemize}

\begin{tcolorbox}[breakable,colback=teal!4,colframe=teal!35!black,title={Before you move on}]
What does o peão mean? (\textquotedblleft{}The pedestrian; the pawn (in chess)\textquotedblright{}.) Say it once more.
\end{tcolorbox}
\section[a passadeira]{a passadeira — the zebra crossing}

\label{lesson:PT-C408-passadeira}



\begin{quote}
Before the new one: say the Portuguese for to unscrew, then the Portuguese for the pedestrian.
\end{quote}

\subsection*{You'll want to know: a passadeira}
\textbf{a passadeira} — \textquotedblleft{}the zebra crossing\textquotedblright{}.

\begin{quote}
O carro parou para nós na passadeira.
\end{quote}

\begin{quote}
\emph{The car stopped for us at the zebra crossing.}
\end{quote}

In Brazil: \emph{a faixa de pedestres}. \emph{Uma passadeira} is also a long, narrow rug.

\begin{cousinweb}
\emph{passadeira} is built on \emph{passar} 'to cross, to pass' (EtymDB) with \emph{-deira}, a place for doing it; \emph{passar} goes back to Latin \emph{passus} 'step'.
\end{cousinweb}

\begin{grammarlens}[title={What you've built}]
One more word for telling, explaining and describing at B1.
\end{grammarlens}

\subsection*{Your turn}
\begin{itemize}
\raggedright
  \item \emph{Say it:} \emph{a passadeira}
  \item \emph{Say it:} \emph{a passadeira}, once more
  \item \emph{Say it:} a sentence of your own with \emph{a passadeira}
  \item \emph{Recall:} say the Portuguese for to unscrew, then the Portuguese for the pedestrian, then say \emph{a passadeira} again
\end{itemize}

\begin{tcolorbox}[breakable,colback=teal!4,colframe=teal!35!black,title={Before you move on}]
What does a passadeira mean? (\textquotedblleft{}The zebra crossing\textquotedblright{}.) Say it once more.
\end{tcolorbox}
\section[atropelar]{atropelar — to run over, to knock down}

\label{lesson:PT-C408-atropelar}



\begin{quote}
Before the new one: say the Portuguese for the pedestrian, then the Portuguese for the zebra crossing.
\end{quote}

\subsection*{You'll want to know: atropelar}
\textbf{atropelar} — \textquotedblleft{}to run over, to knock down\textquotedblright{}.

\begin{quote}
Um carro quase atropelou um cão na passadeira.
\end{quote}

\begin{quote}
\emph{A car almost ran over a dog on the zebra crossing.}
\end{quote}

\emph{Ser atropelado} — to be run over.

\begin{cousinweb}
\emph{atropelar} is built from \emph{a-} and \emph{tropel} 'rush, trampling crowd'. The origin of \emph{tropel} is not settled; EtymDB has no entry for these words.
\end{cousinweb}

\begin{grammarlens}[title={What you've built}]
One more word for telling, explaining and describing at B1.
\end{grammarlens}

\subsection*{Your turn}
\begin{itemize}
\raggedright
  \item \emph{Say it:} \emph{atropelar}
  \item \emph{Say it:} \emph{atropelar}, once more
  \item \emph{Say it:} a sentence of your own with \emph{atropelar}
  \item \emph{Recall:} say the Portuguese for the pedestrian, then the Portuguese for the zebra crossing, then say \emph{atropelar} again
\end{itemize}

\begin{tcolorbox}[breakable,colback=teal!4,colframe=teal!35!black,title={Before you move on}]
What does atropelar mean? (\textquotedblleft{}To run over, to knock down\textquotedblright{}.) Say it once more.
\end{tcolorbox}
\section[o radar]{o radar — the speed camera; the radar}

\label{lesson:PT-C408-radar}



\begin{quote}
Before the new one: say the Portuguese for the zebra crossing, then the Portuguese for to run over.
\end{quote}

\subsection*{You'll want to know: o radar}
\textbf{o radar} — \textquotedblleft{}the speed camera; the radar\textquotedblright{}.

\begin{quote}
Há um radar à entrada da cidade.
\end{quote}

\begin{quote}
\emph{There is a speed camera at the entrance to the town.}
\end{quote}

Also the radar of a ship or an airport.

\begin{cousinweb}
\emph{radar} was borrowed from English \emph{radar}, a word made from the first letters of \emph{radio detection and ranging} (EtymDB lists English \emph{radar}).
\end{cousinweb}

\begin{grammarlens}[title={What you've built}]
One more word for telling, explaining and describing at B1.
\end{grammarlens}

\subsection*{Your turn}
\begin{itemize}
\raggedright
  \item \emph{Say it:} \emph{o radar}
  \item \emph{Say it:} \emph{o radar}, once more
  \item \emph{Say it:} a sentence of your own with \emph{o radar}
  \item \emph{Recall:} say the Portuguese for the zebra crossing, then the Portuguese for to run over, then say \emph{o radar} again
\end{itemize}

\begin{tcolorbox}[breakable,colback=teal!4,colframe=teal!35!black,title={Before you move on}]
What does o radar mean? (\textquotedblleft{}The speed camera; the radar\textquotedblright{}.) Say it once more.
\end{tcolorbox}
\section[fiscalizar]{fiscalizar — to inspect, to check (officially)}

\label{lesson:PT-C408-fiscalizar}



\begin{quote}
Before the new one: say the Portuguese for to run over, then the Portuguese for the speed camera.
\end{quote}

\subsection*{You'll want to know: fiscalizar}
\textbf{fiscalizar} — \textquotedblleft{}to inspect, to check (officially)\textquotedblright{}.

\begin{quote}
A polícia fiscaliza a velocidade nesta estrada.
\end{quote}

\begin{quote}
\emph{The police check speeds on this road.}
\end{quote}

\emph{A fiscalização} — the inspection, the checks.

\begin{cousinweb}
\emph{fiscalizar} is built on \emph{fiscal} 'inspector', from \emph{fisco} 'treasury', which goes back to Latin \emph{fiscus} 'basket; treasury' (EtymDB).
\end{cousinweb}

\begin{grammarlens}[title={What you've built}]
One more word for telling, explaining and describing at B1.
\end{grammarlens}

\subsection*{Your turn}
\begin{itemize}
\raggedright
  \item \emph{Say it:} \emph{fiscalizar}
  \item \emph{Say it:} \emph{fiscalizar}, once more
  \item \emph{Say it:} a sentence of your own with \emph{fiscalizar}
  \item \emph{Recall:} say the Portuguese for to run over, then the Portuguese for the speed camera, then say \emph{fiscalizar} again
\end{itemize}

\begin{tcolorbox}[breakable,colback=teal!4,colframe=teal!35!black,title={Before you move on}]
What does fiscalizar mean? (\textquotedblleft{}To inspect, to check (officially)\textquotedblright{}.) Say it once more.
\end{tcolorbox}
\section[(dialogue)]{Estradas e estações}

\label{lesson:PT-R408-revisao-401-408}



\begin{quote}
Say the Portuguese for the lane and the hard shoulder.
\end{quote}

\subsection*{Your turn}
Say these aloud:

\begin{itemize}
\raggedright
  \item the lane, the hard shoulder, the central reservation, the junction, to speed up
  \item the brake, to skid, the pedal, the clutch, to put into gear
  \item the glass, the rear-view mirror, the horn, to honk, to mist up
  \item diesel, fuel, the hose, to fill up, to fill to the top
  \item the tyre, to puncture, the puncture, the wheel rim, the warning triangle
  \item the breakdown, to break down, to tow, the tow truck, the repair garage
  \item the bonnet, the boot, the indicator, the screw, to unscrew
  \item the pedestrian, the zebra crossing, to run over, the speed camera, to inspect
\end{itemize}

\begin{tcolorbox}[breakable,colback=teal!4,colframe=teal!35!black,title={Before you move on}]
The lane? (\textbf{a faixa}.) The hard shoulder? (\textbf{a berma}.) The central reservation? (\textbf{o separador}.) The junction? (\textbf{o nó}.) To speed up? (\textbf{acelerar}.) The brake? (\textbf{o travão}.) To skid? (\textbf{derrapar}.) The pedal? (\textbf{o pedal}.) The clutch? (\textbf{a embraiagem}.) To put into gear? (\textbf{engatar}.) The glass? (\textbf{o vidro}.) The rear-view mirror? (\textbf{o retrovisor}.) The horn? (\textbf{a buzina}.) To honk? (\textbf{buzinar}.) To mist up? (\textbf{embaciar}.) Diesel? (\textbf{o gasóleo}.) Fuel? (\textbf{o combustível}.) The hose? (\textbf{a mangueira}.) To fill up? (\textbf{abastecer}.) To fill to the top? (\textbf{atestar}.) The tyre? (\textbf{o pneu}.) To puncture? (\textbf{furar}.) The puncture? (\textbf{o furo}.) The wheel rim? (\textbf{a jante}.) The warning triangle? (\textbf{o triângulo}.) The breakdown? (\textbf{a avaria}.) To break down? (\textbf{avariar}.) To tow? (\textbf{rebocar}.) The tow truck? (\textbf{o reboque}.) The repair garage? (\textbf{a oficina}.) The bonnet? (\textbf{o capô}.) The boot? (\textbf{a bagageira}.) The indicator? (\textbf{o pisca}.) The screw? (\textbf{o parafuso}.) To unscrew? (\textbf{desaparafusar}.) The pedestrian? (\textbf{o peão}.) The zebra crossing? (\textbf{a passadeira}.) To run over? (\textbf{atropelar}.) The speed camera? (\textbf{o radar}.) To inspect? (\textbf{fiscalizar}.)
\end{tcolorbox}
